\documentclass[10pt,a4paper]{amsart}
\usepackage[margin=19mm]{geometry}
\usepackage{amsmath,amssymb,mathtools}
\usepackage[scr=rsfs,cal=euler]{mathalfa}
\usepackage{xfrac}
\usepackage[unicode,hidelinks,bookmarks=false]{hyperref}
\newcommand{\ie}{i.e.\ }
\newcommand{\cf}{cf.\ }
\newcommand{\resp}{respectively\ }
\newcommand{\Subsection}{\subsection}
\providecommand{\nobreakdash}{\nobreak\hbox{-}}
\setlength{\emergencystretch}{2em}
\multlinegap0pt
\usepackage{bm}
\usepackage{bbm}
\usepackage{dsfont}
\usepackage{xspace}
\usepackage{cancel}
\usepackage{mathtools}
\usepackage{cleveref}
\usepackage{amsthm}
\makeatletter
\@ifundefined{theorem}{\newtheorem{theorem}{Theorem}}{}
\@ifundefined{lemma}{\newtheorem{lemma}[theorem]{Lemma}}{}
\@ifundefined{proposition}{\newtheorem{proposition}[theorem]{Proposition}}{}
\makeatother
\makeatletter
\expandafter\ifx\csname cenumeratei\endcsname\relax
\newenvironment{cenumeratei}%
    {\begin{list}{\labelcenumi}{\usecounter{cenumi}\partopsep=4pt\topsep=4pt\itemsep=4pt\parsep=4pt}%
    \setcounter{cenumi}{\value{cenumisaved}}}%
      {\setcounter{cenumisaved}{\value{cenumi}}%
    \end{list}}
\fi
\expandafter\ifx\csname cenumerateii\endcsname\relax
\newenvironment{cenumerateii}%
    {\begin{list}{\labelcenumii}{\usecounter{cenumii}\partopsep=0pt\topsep=0pt\itemsep=0pt\parsep=0pt}%
    }%
      {%
    \end{list}}
\fi
\expandafter\ifx\csname proof\endcsname\relax
\renewenvironment{proof}{\par\noindent{\bf Proof\ }}{\hfill\BlackBox\\[2mm]}
\fi
\makeatother
\title[Associative Memory and Generative Diffusion in the Zero-nois]{Associative Memory and Generative Diffusion in the Zero-noise Limit}
\author{Joshua Hess, Quaid Morris}
\date{Source: arXiv:2506.05178v2 (CC BY 4.0)}
\begin{document}
\maketitle

\noindent\emph{Source and licence.} Associative Memory and Generative Diffusion in the Zero-noise Limit. Version arXiv:2506.05178v2 (CC BY 4.0), distributed under the Creative Commons Attribution 4.0 International licence, as recorded in the arXiv licence metadata for this exact version and shown on the abstract page https://arxiv.org/abs/2506.05178v2. Full rights notice: licenses/arxiv-2506.05178-CC-BY-4.0.txt.\par\smallskip

\noindent\textit{Editorial scope.} Counted unit: the Proposition whose source label is \texttt{lemma:zero-noise-continuous-dependence} in the article (source0.tex lines 733--737, source label \texttt{lemma:zero-noise-continuous-dependence}), delivered together with the article's own complete original proof of it (source0.tex lines 739--759, one \texttt{proof} environment of 21 lines). No mathematics is written, repaired or re-derived: statement and proof are the article's own text line for line. Required context retained uncounted: equation:vector-field-ode-definition (lines 483--486); theorem:stable-manifold-theorem-morse (lines 603--611); proposition:zero-noise-limit-MS (lines 658--661); lemma:good-critical-points-milnor (lines 1575--1578). Every definition, display and label of those lines is retained unchanged. Omitted from this package and not redistributed: 1--482, 487--602, 612--657, 662--732, 738--738, 760--1574, 1579--1748. Nothing inside a retained excerpt is omitted or altered: every retained body is delivered exactly as the article's lines are written, so no omission receipt of any kind accompanies this package. Citations: the retained excerpts carry no citation command at all, so no bibliography entry of the article is transcribed and no reference is redistributed. Reference adaptation: none was needed and none was made; every reference inside the delivered unit resolves inside it, so no unresolved reference or citation is produced. Printed slips kept literal: no printed word, symbol, hypothesis or number of the counted statement or of its proof was altered. Layout and class: the article's own class is \texttt{article} with packages including amsmath, times, graphicx, color, multirow, natbib, rotating, bbm, latexsym, mathrsfs, mathtools, tikz, subfig, setspace; the wrapper is the lane's pinned \texttt{amsart} editorial wrapper at 10pt on A4, which restates the theorem-like environments the retained text needs and adds the article's own macro definitions that the retained text needs (none), with extra wrapper packages (bm, bbm, dsfont, xspace, cancel, mathtools, cleveref). The delivered numbering is the unit's own. Assets: no retained span contains a figure environment, a graphics-inclusion command, a PDF-page inclusion, a table or a TikZ picture; no image file is copied into this package, the package declares no asset dependency and no image file is redistributed with it. Licence: version arXiv:2506.05178v2 of this article is distributed under the Creative Commons Attribution 4.0 International licence, as recorded in the arXiv licence metadata for this exact version and shown on the abstract page https://arxiv.org/abs/2506.05178v2. A full rights notice is delivered at licenses/arxiv-2506.05178-CC-BY-4.0.txt. Characters: every retained excerpt is pure ASCII, so the delivered engine, XeLaTeX with its default Latin Modern text font, reports no missing character.
\begin{equation}
\label{equation:vector-field-ode-definition}
    \frac{d}{dt}\phi^X_t(p) = X\left(\phi^X_t(p) \right), \hspace{0.5em}  \text{with }\phi^X_0(p)=p \,.
\end{equation}

\begin{theorem}[Stable Manifold Theorem for Morse functions]
\label{theorem:stable-manifold-theorem-morse}
    Let $X = -\nabla_g V$ be a gradient field on a Riemannian $n$-manifold $(M,g)$ with $V \in C^{r+1}(M,\mathbb{R}), r \geq 1$ a Morse function. Let $p \in M$ be a hyperbolic fixed point of $V$ with index $\lambda_p$, and $E^s$ (resp. $E^u$) the stable (resp. unstable) subspace of $DX_p = L$. Then the following hold,
    \begin{enumerate}
        \item $W^s(p)$ is an embedded differentiable manifold in $M$ with the same regularity as $\phi^X_t$, and the tangent space $T_p(W^s(p))$ at $p$ is $E^s$; hence, 
        $W^s(p)$ and $W^u(p)$ are embedded (open) topological discs with dimensions $n-\lambda_p$ and $\lambda_p$, respectively.
        \item Let $D \subset W^s(p)$ be an embedded disc containing $p$. Then there is a neighborhood $\mathcal{N} \subset \text{Diff}^r(M)$ in which all $g \in \mathcal{N}$ have a unique hyperbolic fixed point $p_g$ contained in a neighborhood $U \ni p$. Hence, for any $\epsilon > 0$, there is a neighborhood $\tilde{\mathcal{N}} \subset \mathcal{N}$ of $\phi_t^X$ such that, for each $g \in \tilde{\mathcal{N}}$, there exists a disc $D_g \subset W^s(p_g)$ that is $\epsilon$-$C^r$ close to $D$.
    \end{enumerate}
\end{theorem}

\begin{proposition}[Generic zero-noise limits.]
\label{proposition:zero-noise-limit-MS}
    Let $(M,g)$ be a Riemannian manifold and $X = -\nabla_g V$ a gradient vector field on $M$ with $V\in C^{r}(M,\mathbb{R})$ a Morse function where $r \geq 2$; let $\{\mathcal{X}^{\epsilon}\}_{\epsilon>0}$ be a family of small random perturbations of $\phi_t^X$ with invariant measures $\mu^{\epsilon}$ for $\mathcal{X^{\epsilon}}$ at a given $\epsilon>0$. Then as $\epsilon \rightarrow 0$, $\{\mu^{\epsilon}\}_{\epsilon>0}$ converges weakly to an invariant measure $\mu$ of $\phi_t^X$ consisting of a weighted sum $\mu=\sum_i^n w_i\delta_{\beta_i}$, where $\{ \beta_1,...,\beta_n \} = \Omega(X)$ is the non-wandering set of $X$.
\end{proposition}

\begin{lemma}
\label{lemma:good-critical-points-milnor}
    Let $K$ be a compact subset of an open set $U \subset \mathbb{R}^n$ and $1 \leq r \leq \infty$. If $f\in C^r(U,\mathbb{R})$ has only nondegenerate critical points in $K$, then there exists an $\eta >0$ such that if $g\in C^r(U,\mathbb{R})$ and $|| \, D^i(f-g)(x) \, || < \epsilon$ for all $x\in K$, then $g$ has only nondegenerate critical points in $K$.
\end{lemma}

\begin{proposition}
[Continuous dependence of zero-noise limits]
\label{lemma:zero-noise-continuous-dependence}
    Let $\phi^X_t$ be a Morse gradient flow generated by $X = -\nabla_g V$ for $V \in C^{r+1}(M,\mathbb{R}), r\geq1$ a Morse function. Denote the set of critical elements of $X$ by $\{ \beta_i,...,\beta_n\}$ and let $W^s(\beta_a)$ be a stable manifold for $\beta_a$, $a \in 1,...,n$. If $\{\mathcal{X}^{\epsilon}\}_{\epsilon>0}$ is a family of small random perturbations of $\phi_t^X$. Then the zero noise limit of $\{\mathcal{X}^{\epsilon}\}_{\epsilon>0}$ on $W^s(\beta_a)$ depends continuously on the $C^r$ flow $\phi^X_t$ for the weak topology on measures and the compact-open $C^{r+1}$ topology on real-valued functions.
\end{proposition}

\begin{proof}
    Let $\mu$, and $\mu'$ be zero-noise limits for small random perturbations $\{ \mathcal{X^{\epsilon}}\}_{\epsilon >0}$ and $\{ \mathcal{X'^{\epsilon}}\}_{\epsilon >0}$ of the gradient flows $\phi_t^X$ and $\phi_t^{X'}$, respectively, with $X' \in \text{Grad}^r(M)$. We have to show that the sequence $\{\mu'\}$ converges to $\mu$ weakly on $W^s(\beta_a)$ when $V' \rightarrow V$ in the $C^{r+1}$ sense, where $X' = -\nabla_g V'$. Let $(U_1,\psi_1),...,(U_k,\psi_k)$ be a finite covering of $M$ (by compactness). Choose compact sets $K_i \subset U_i$ such that $\bigcup_i K_i$ covers $M$. From \Cref{lemma:good-critical-points-milnor}, there exists an $\alpha>0$ so that all 
    functions $V' \in C^{r+1}(M,\mathbb{R})$ with
    \begin{equation*}
        || D^j (V \circ \psi^{-1})(x) - D^j(V' \circ \psi^{-1})(x)|| < \alpha
    \end{equation*}
    have nondegenerate critical elements on $M$, where $x \in \psi(K)$ with $K = K_i$ for some $i$ and $j = 0,...,r+1$. This $\alpha$-neighborhood of $V$ induces a neighborhood of $\nabla_gV$ in $C^r(M,TM)$ and similarly on flows by the uniqueness of the solutions to \eqref{equation:vector-field-ode-definition}. That is, for $V'$ within $\alpha$ of $V$ in the $C^{r+1}$ sense above, the flows $\phi_t^X$ and $\phi_t^{X'}$ are within some $\eta >0$. For any $\delta \leq \eta$, non-degeneracy of critical points is preserved since $V'$ is within $\alpha$ of $V$. Call this neighborhood $\mathcal{U}$. Convergence in the compact-open $C^r$ topology is uniform. Therefore for any $\delta>0$, there exists an $N \in \mathbb{Z}^{+}$, such that for all $n \geq N$,
    \begin{equation*}
        || D^l (\psi \circ \phi_t^{X'_{n}} \circ \psi^{-1})(x) - D^l (\psi \circ \phi_t^{X} \circ \psi^{-1})(x) || < \delta \, ,
    \end{equation*}
    for $x \in \psi(K_i)$ for some $K_i$ and $l=0,...,r$. Call this $\delta$-neighborhood $\mathcal{V}$ and let $\mathcal{N = \mathcal{U} \, \cap \, \mathcal{V}}$. Then for any $\phi_t^{X'_{n}}\in \mathcal{N}$, the stable manifolds of its critical elements are well-defined, since the set $\text{Crit}(V'_{n})$ contains no degenerate critical points. By \Cref{proposition:zero-noise-limit-MS}, the zero-noise limit $\mu_n'$ is a weighted sum of Dirac delta functions. By the Stable Manifold Theorem for a Morse function (\Cref{theorem:stable-manifold-theorem-morse}), $\delta$ can be chosen sufficiently small so the stable manifold $W^s(\beta_a^{'n})$ is $\epsilon$-$C^r$ close to $W^s(\beta_a)$ for any $\epsilon > 0$. That is,
    $W^s(\beta_a^{'n}) \rightarrow W^s(\beta_a)$ in the $C^r$ sense as $V'_n \rightarrow V$ in the $C^{r+1}$ sense, and $\beta_a^{'n} \rightarrow \beta_a$. On $W^s(\beta_a)$, $\mu$ is written $\mu = \int_{W^s(\beta_a)} w_a \delta_{\beta_{a}} \left(d\text{vol}_g(x)\right)$ subject to the constraint $\mu(W^s(\beta_a))=1$. Therefore, $w_a=1$. The result follows immediately since $\beta_a^{'n} \rightarrow \beta_a$ and by \Cref{proposition:zero-noise-limit-MS}, $\mu'_n$ is similarly written $\mu'_n = \int_{W^s(\beta^{'n}_a)} w^{'n}_a \delta_{\beta^{'n}_{a}} \left(d\text{vol}_g(x)\right)$ with $w'_a=1$. We have that for any $f \in C^0_b(M)$,
    \begin{equation*}
        \int_{W^s(\beta_a^{'n})} f \,d\mu'_{n} = w_a^{'n} f(\beta_a^{'n}) \rightarrow w_a f(\beta_a) = \int_{W^s(\beta_a)} f \,d\mu.
    \end{equation*}
    Therefore, the subsequence $\{\mu'_{n}\}$ converges weakly to $\mu$. To show that $\{\mu'\}$ converges weakly to $\mu$, proceed by contradiction. Suppose that $\{\mu'\}$ does not converge to $\mu$. Then there exists a $g\in C^0_b(W^s(\beta_a))$, a number $\delta > 0$, and a subsequence $\{\mu'_{k}\}$ so that 
    \begin{equation*}
        \left|\int_{W^s(\beta_a)} g \,d\mu - \int_{W^s(\beta_a^{'k})} g \,d\mu'_{k} \right| \geq \delta
    \end{equation*}
    for all $k \geq 1$. Therefore no subsequence converges, which is a contradiction.
\end{proof}

\end{document}
